\chapter{Night Terrors (Sleep Terrors)}

\section{Definition}
Night terrors are a parasomnia disorder characterized by abrupt arousal from deep non-REM sleep with intense fear, screaming, and autonomic activation. The person is inconsolable and has no memory of the event upon waking. Night terrors most commonly affect children aged 3-7 years.

\section{Pathophysiology}
Night terrors occur during arousal from slow-wave (stage N3) sleep, typically in the first third of the night. There is a partial arousal from deep sleep with activation of the sympathetic nervous system (tachycardia, tachypnea, sweating, mydriasis). The brain remains in a mixed sleep-wake state, with the person unable to fully awaken. Genetic factors play a role.

\section{Etiology}
\begin{itemize}
\item Inadequate sleep or sleep deprivation
\item Fever or illness
\item Stress or anxiety
\item Irregular sleep schedule
\item Sleeping in an unfamiliar environment
\item Bladder distension
\item Genetics: family history of parasomnias
\item Underlying sleep disorders (obstructive sleep apnea, restless legs syndrome)
\item Febrile illness
\item Medications (sedatives, hypnotics, psychostimulants)
\end{itemize}

\section{Indications for Evaluation}
\begin{itemize}
\item Sudden sitting up in bed with a frightened expression
\item Loud screaming or crying during sleep
\item Intense autonomic arousal: rapid heart rate, sweating, dilated pupils
\item Inconsolable during the episode
\item No response to comforting or verbal communication
\item Episode lasts 1-10 minutes
\item Returns to sleep spontaneously without full awakening
\item Complete amnesia for the event in the morning
\end{itemize}

\section{Related Symptoms}
\begin{itemize}
\item Sleepwalking (somnambulism)
\item Sleep talking
\item Confusional arousals
\item Nightmare disorder
\item Sleep apnea
\item Restless legs syndrome
\end{itemize}

\section{Self-Care/Home Management}
\begin{itemize}
\item Ensure adequate sleep (children need 10-12 hours)
\item Maintain a consistent bedtime routine
\item Create a calm, relaxing sleep environment
\item Do not wake the child during a night terror (prolongs the episode)
\item Gently guide the child back to bed and stay nearby
\item Protect the child from injury (lock doors, remove obstacles)
\item Reduce stress and screen time before bed
\end{itemize}

\section{Diagnostic Approach}
\begin{itemize}
\item Clinical history from parents or bed partner
\item Polysomnography for atypical, frequent, or injurious episodes
\item Differentiate from nightmares (which occur in REM sleep and are recalled)
\item EEG to rule out nocturnal seizures
\end{itemize}

\section{Treatment/Management Overview}
\begin{itemize}
\item Reassurance and education for parents (spontaneous resolution is common)
\item Address underlying sleep deprivation and schedule irregularity
\item Treat coexisting sleep disorders
\item Scheduled awakening (waking the child 15-30 minutes before typical episode time)
\item Benzodiazepines (clonazepam) or tricyclic antidepressants for severe, frequent episodes
\end{itemize}
